\chapter{Economic, Regulatory and Methodological Context}
\label{chap:context}

This chapter sets out the context in which the project takes shape. It first
describes cyber crises as a growing concern for European organisations, then
examines how the NIS2 Directive shapes the regulatory expectations that apply
to crisis preparedness, and how the ENISA \emph{Best Practices for Cyber
Crisis Management} framework responds to that expectation. From there, the
chapter identifies the gap between what this framework prescribes and what
organisations can currently measure, and closes with the precise problem
statement addressed by this thesis. Existing tools and research addressing
comparable problems are reviewed separately in \ref{chap:state-of-the-art}.

\section{The cyber crisis landscape}
\label{sec:crisis-landscape}

\subsection{A growing concern}

Cyber incidents have moved from an operational nuisance to a standing concern
for European organisations of all sizes.
% TODO: optionally add a citable figure here if you have one you can verify
%   directly at the source, the paragraph reads fine without one (Ch.3's
%   verification notes dropped a similar unconfirmed statistic rather than
%   force one in, see REF11 in verification-references-chapitre3.md).
Across the sector, the question organisations ask is shifting accordingly,
from whether a serious incident will occur to how well they are able to
respond when it does.

\subsection{Cybersecurity and crisis management: two distinct capabilities}

This shift in framing separates two related but distinct concerns.
Cybersecurity, in its usual sense, is largely about reducing the likelihood
and the technical impact of an incident, through system hardening, intrusion
monitoring and vulnerability management. Crisis management is a different
capability. It is the organisation's ability to make timely decisions with
incomplete information, coordinate stakeholders across technical,
operational and strategic levels, and communicate under pressure, once an
incident has escalated into an event that threatens the organisation's
ability to operate. A well-secured organisation can still handle a crisis
poorly, and a mature crisis management capacity can limit the damage of an
incident that prevention failed to stop.

This distinction is more than a semantic one. The referential underpinning
this project, presented in detail in \ref{chap:scoring-engine}, treats two
dimensions as cutting across all four phases of crisis management:
decision-making under uncertainty, and an all-hazards approach that accounts
for physical, third-party and human-failure scenarios alongside cyber-native
threats. Both point toward the same conclusion. Crisis management calls for
a maturity assessment distinct from a general cybersecurity posture
assessment, one that this thesis sets out to build.

Increasing regulatory attention reflects this same shift. The following
section examines how the NIS2 Directive translates this concern into a set
of operational obligations.

\section{The NIS2 Directive and its operational implications}
\label{sec:nis2}

\subsection{Scope and obligations}

The pressure examined in \ref{sec:crisis-landscape} increasingly takes a
regulatory form. The NIS2 Directive extends the scope and depth of
cybersecurity obligations for organisations operating in sectors deemed
essential or important to the Union's economy and society
\cite{nis2_directive_2022}.

Among the obligations NIS2 introduces, crisis management sits alongside more
familiar requirements such as incident handling and business continuity
\cite{nis2_directive_2022}.
% TODO(fact): optional. If you want to cite a specific article/paragraph
%   (commonly associated with Article 21) rather than the directive as a
%   whole, confirm the exact reference first. Chapter 3 only cites the
%   directive generically for comparable claims, so this level of detail is
%   not required to stay consistent with the rest of the thesis.
This is what gives the project its regulatory anchor. Crisis management is
not an optional maturity dimension an organisation may choose to develop on
its own schedule, it is one NIS2 asks regulated entities to be able to
demonstrate.

\subsection{What demonstrating maturity requires}

Demonstrating this capability calls for more than a policy document kept on
a shelf. Regulators and auditors increasingly expect evidence, documented
processes, tested plans, traceable exercises, rather than declarative
claims of readiness. Entities are told what they are expected to
demonstrate, without being given a way to measure how close they currently
are to that standard.

The next section examines how the ENISA \emph{Best Practices for Cyber
Crisis Management} framework addresses what organisations are expected to
put in place \cite{enisa_crisis_mgmt_2024}. Section~\ref{sec:gap} then turns
to what is still missing, a way to measure where an organisation currently
stands against it.

\section{The ENISA \emph{Best Practices for Cyber Crisis Management} framework}
\label{sec:enisa-framework}

\subsection{Structure of the framework}

The regulatory pressure examined in \ref{sec:nis2} calls for an operational
reference specific to crisis management, rather than cybersecurity in
general. ENISA answered this need in February 2024 with the \emph{Best
Practices for Cyber Crisis Management} study \cite{enisa_crisis_mgmt_2024},
a document aimed squarely at the discipline described in
\ref{sec:crisis-landscape}.

The framework organises cyber crisis management into four phases,
Prevention, Preparedness, Response and Recovery, each covering a distinct
moment in the life of a crisis rather than a distinct technical domain.
Figure~\ref{fig:enisa-phases} summarises this structure as encoded in the
referential built for this project (its full construction is detailed in
\ref{chap:scoring-engine}): each phase carries an equal 25\% weight in the
resulting maturity score, regardless of how many best practices it
contains. Sixteen best practices are distributed unevenly across the four
phases, from three in Prevention to nine in Preparedness, reflecting the
fact that most of what determines crisis readiness is decided before a
crisis starts, not during it.

\begin{figure}[ht!]
  \centering
  \begin{tikzpicture}[
    phase/.style={
      draw, thick, rounded corners, minimum width=3.1cm, minimum height=1.5cm,
      align=center, font=\sffamily\small
    },
    lbl/.style={font=\sffamily\scriptsize, align=center},
    arr/.style={-{Stealth[length=2.5mm]}, thick}
  ]

  \node[phase] (prev) at (0, 3.6) {Prevention};
  \node[phase] (prep) at (3.6, 0) {Preparedness};
  \node[phase] (resp) at (0, -3.6) {Response};
  \node[phase] (recv) at (-3.6, 0) {Recovery};

  \node[lbl, above=1mm of prev] {25\% (3 BPs)};
  \node[lbl, right=1mm of prep] {25\% (9 BPs)};
  \node[lbl, below=1mm of resp] {25\% (2 BPs)};
  \node[lbl, left=1mm of recv] {25\% (2 BPs)};

  \draw[arr] (prev) to[bend left=18] (prep);
  \draw[arr] (prep) to[bend left=18] (resp);
  \draw[arr] (resp) to[bend left=18] (recv);
  \draw[arr] (recv) to[bend left=18] (prev);

  \node[lbl, font=\sffamily\footnotesize\itshape] at (0,0)
    {Cyber crisis\\management cycle};

  \end{tikzpicture}
  \caption{The four phases of the ENISA cyber crisis management framework, each normalised to an equal 25\% weight in the global maturity score regardless of the number of best practices it contains (source: referential.yaml, derived from \cite{enisa_crisis_mgmt_2024})}
  \label{fig:enisa-phases}
\end{figure}

\subsection{Foundational best practices}

Within this structure, four best practices are treated as foundational,
BP04, BP07, BP08 and BP12, covering respectively governance and the crisis
coordinator role, allocation of roles and responsibilities, escalation
criteria, and communication strategy. The referential built for this
project enforces this distinction directly: an organisation cannot be rated
highly mature overall if any of these four best practices scores too low,
regardless of how well the remaining twelve perform (the mechanism is
detailed in \ref{chap:scoring-engine}).
% TODO(fact): if you want to explicitly name or contrast NIS2's own
%   national-level coordination provisions here, confirm the exact article
%   first (contexte-memoire flags Art. 9 as a possibility, unconfirmed). The
%   sentence below deliberately stays at a level that does not require that
%   confirmation.
The crisis coordinator introduced by BP04 is an internal governance role
within the organisation being assessed, distinct from any coordination
mechanism NIS2 may establish at Member State level.

Framed this way, the ENISA framework is thorough and directly relevant to
the regulatory context described above, but it remains descriptive. It
specifies which best practices an organisation should have in place,
without offering a scale, an evidence standard, or a way to turn sixteen
qualitative assessments into a single defensible score. This is the gap the
next section describes in detail.

\section{The gap addressed by this project}
\label{sec:gap}

\subsection{A specific, identifiable gap}

The previous section closed on a specific limitation: ENISA specifies what
good crisis management should look like, without specifying how to measure
it. Generic cybersecurity maturity tools do not fill this limitation
either, they treat crisis management as one dimension among many rather
than as a discipline in its own right, and none of them maps natively to
the ENISA framework or to NIS2, a point examined in detail in
\ref{chap:state-of-the-art}. What remains, once the existing landscape is
set aside, is a specific and identifiable gap: an organisation, or a
consultant working with one, that wants to assess crisis-management
maturity against the reference EU-level framework has no purpose-built
instrument to do so.

\subsection{Why it matters, for consultants and clients alike}

This gap matters from two converging angles. On the consulting side, a
crisis-readiness conversation with a client currently opens with generic
checklists or with a consultant's own judgement, neither of which produces
a result the client can later audit or compare across engagements. What is
missing is a quick, defensible starting point, quick enough to fit an
early-stage engagement, defensible enough to withstand scrutiny of how the
result was produced.

On the client side, the regulatory pressure examined in \ref{sec:nis2}
creates a mirrored need. An organisation that must be able to demonstrate
crisis-management maturity requires something more structured than a
self-declared statement, without the cost and lead time of a full external
audit at every stage of its NIS2 compliance journey.

Both angles converge on the same requirement, a way to turn the sixteen
ENISA best practices into an evidence-grounded, auditable score. The
following section states this requirement as the precise problem this
thesis addresses.

\section{Problem statement}
\label{sec:problem-statement}

The chapter can now be summarised into a single problem. Organisations
subject to NIS2 are expected to demonstrate crisis-management maturity
against the ENISA \emph{Best Practices for Cyber Crisis Management}
framework (\ref{sec:enisa-framework}), but that framework describes what
good practice looks like without providing a way to measure it
(\ref{sec:gap}). Consultants and client organisations alike therefore need
an instrument that turns the framework's sixteen best practices into a
maturity score that is fast to produce, grounded in evidence rather than
declarative claims, and able to hold up when a client's own governance and
security leadership examine how it was produced.

Three constraints follow directly from this framing, and shape everything
the rest of this thesis builds. Regulatory pressure means the result must
survive scrutiny, not just sound plausible in conversation: a client's
governance and security leadership, not only the consultant presenting it,
must be able to trace every score back to the evidence and the rule that
produced it (\ref{sec:nis2}). Technical fidelity means the instrument must
stay grounded in the ENISA framework itself rather than in an ad hoc scale
invented for convenience, the referential, not the tool, has to remain the
source of truth (\ref{sec:enisa-framework}). Methodological accessibility
means the instrument must be usable by a consultant leading a client
conversation, not only by someone with deep prior expertise in the ENISA
text, which rules out an assessment that requires extensive training to
administer or interpret.

Meeting these constraints requires, at minimum, three outcomes: a scoring
method that is deterministic and repeatable across assessments, a bridge
that maps assessment answers into the framework without hardcoding its
content into whatever tool implements it, and a way to surface the
resulting score and its underlying evidence to both the person answering
the assessment and the person relying on its result. How these outcomes are
achieved, through which principles, architecture and implementation
choices, is the subject of \ref{chap:method}, \ref{chap:scoring-engine} and
\ref{chap:platform}.

Before turning to how this problem is addressed, it is worth situating it
against what already exists. Existing maturity tools, crisis-management
research and the constraints of LLM-based augmentation are reviewed in
\ref{chap:state-of-the-art}.